% ============================================================
% sec/06etica.tex — VERSIÓN INTEGRADA, TRES EJES
%
% OBLIGATORIO: riesgos identificados y forma prevista de
% mitigarlos. La detección corresponde al momento previo al
% entrenamiento.
% ============================================================
\section{Consideraciones Éticas Preliminares}
\label{sec:etica}

\subsection{Sensibilidad de los datos}
\label{subsec:sensibilidad}

El corpus no contiene datos clínicos, biométricos ni conductuales sobre
individuos identificados: sus registros son partidas simuladas entre agentes
de software, publicadas bajo una licencia CC0-1.0
\cite{ptcg_dataset_2026}. En la auditoría realizada no se identificaron datos
personales sensibles de jugadores humanos.

Cabe una salvedad. Cada episodio lleva los nombres de equipo y de agente de
los competidores, y en la muestra auditada esos nombres son seudónimos
elegidos por sus dueños que en algunos casos se asemejan a nombres
personales. Identifican a participantes de una competencia pública y no a
sujetos de estudio, pero siguen siendo atribuibles a personas.

Se tratan por tanto como metadatos y no como material: se usan para
auditoría, agrupación y la comprobación de juego contra uno mismo, quedan
excluidos del vector de características por la lista negra descrita en la
Sección~\ref{subsec:leakage}, y no se publicará ningún ordenamiento ni
comparación agregada de competidores nombrados.

\razon{Esta decisión coincide con el control de fuga, lo cual es conveniente
pero no es la razón de la decisión. Declararlo así evita presentar una
precaución técnica como si fuera una consideración ética, o viceversa.}

\subsection{Riesgo de sesgo y representatividad}
\label{subsec:sesgo}

El principal riesgo de este proyecto no es el daño a un sujeto de datos: es
la sobreafirmación.

Los episodios se seleccionan a diario por puntuación media de los agentes y
se truncan a un volumen diario fijo \cite{ptcg_dataset_2026}, de modo que el
corpus es la franja alta de una escalera de agentes artificiales y no una
muestra del juego. La auditoría cuantifica la concentración que se sigue de
ello: los diez mazos más frecuentes cubren el 39.94\,\% de las instancias, un
día auditado contiene solo 48 mazos distintos, y el conjunto de mazos del
último día auditado es esencialmente disjunto del del primero. Un modelo
ajustado aquí aprende un conjunto de estrategias estrecho y en movimiento.

La literatura revisada refuerza que el riesgo no es hipotético. Angliss
\textit{et al.} encuentran diferencias entre el rendimiento sobre
composiciones vistas y la generalización a composiciones no vistas
\cite{angliss_vgcbench_2026}. Chitayat \textit{et al.} separan la evaluación
entre versiones conocidas y posteriores \cite{chitayat_meta_2023}. Xenopoulos
\textit{et al.} muestran que la importancia relativa de las características
difiere entre poblaciones de distinto nivel de habilidad, y que un modelo
entrenado sobre una no necesariamente se extiende a otra sin pérdida de
desempeño \cite{xenopoulos_winprob_2022}.

Saravanan y Guzdial orientan deliberadamente su agente a aproximar un nivel
de habilidad promedio, argumentando que el conjunto de estrategias dominante
constituye una medida de popularidad determinada mayoritariamente por
jugadores de habilidad media \cite{saravanan_metadiscovery_2024}. Un corpus
restringido a la franja superior observa, en consecuencia, una población
distinta de aquella que determina el comportamiento agregado del juego.

\subsubsection{Tres lecturas erróneas, cada una fácil de cometer}

\noindent\textbf{Extrapolar a jugadores humanos} carecería de respaldo: no
aparece ninguna partida humana en el corpus, y los dos corpus comparables de
deportes electrónicos revisados declaran exactamente esta limitación para sí
mismos \cite{bialecki_sc2egset_2023,janusz_hearthstone_2017}.

\noindent\textbf{Extrapolar al juego completo} tampoco tendría respaldo,
puesto que un modelo que nunca vio un arquetipo de mazo carece de base para
juzgar posiciones construidas en torno a él. Bertram \textit{et al.} muestran
cómo se ve ese fallo cuando el contenido se codifica por identidad en lugar
de por descripción \cite{bertram_cardrepresentations_2024}.

\noindent\textbf{Tratar la salida como una propiedad estable de una posición}
tampoco tendría respaldo a lo largo del tiempo, porque la población cambia de
forma medible en diez días.

\subsection{Interpretación de la probabilidad}
\label{subsec:interpretacionprob}

Una estimación de probabilidad de victoria invita a un mal uso específico:
leer un número entre cero y uno como una posibilidad bien fundada cuando
puede ser solo una puntuación bien ordenada.

La distinción no es pedante. Guo \textit{et al.} encontraron deterioro de la
calibración en todos los estudios que la midieron bajo cambio temporal,
frente a deterioro de la discriminación en solo tres de doce
\cite{guo_temporalshift_2021}, lo que significa que un modelo puede seguir
ordenando posiciones correctamente mucho después de que sus probabilidades
declaradas hayan dejado de ser ciertas. La mayoría de los estudios
comparables sobre estados de juego no reporta calibración en absoluto
\cite{janusz_hearthstone_2017,bialecki_sc2egset_2023}.

\razon{Reportar una regla de puntuación propia y curvas de fiabilidad junto
al área bajo la curva es por tanto un compromiso ético tanto como
metodológico: es lo que permite a un lector saber cuál de las dos cosas es
nuestro número.}

\subsection{Límites de generalización declarados}
\label{subsec:limitesgeneralizacion}

Los resultados de este trabajo no deben extrapolarse a:

\begin{itemize}
  \item Partidas entre personas. El corpus procede de agentes artificiales,
        cuyo comportamiento no se asume equivalente al humano.
  \item El \emph{Pok\'emon TCG} en su formato físico.
  \item El conjunto de estrategias posibles del juego, más allá de las
        composiciones observadas.
  \item Poblaciones de agentes con nivel de puntuación distinto del
        observado.
  \item Versiones del motor de simulación distintas de las presentes en los
        datos.
  \item Periodos posteriores a los cubiertos por la muestra auditada, más
        allá de lo que el protocolo B permita evaluar.
  \item Los 57 días no auditados del corpus publicado.
\end{itemize}

Dos precisiones adicionales sobre el alcance de las afirmaciones. Un buen
resultado bajo la partición agrupada no implica robustez temporal: esa
afirmación solo podrá sostenerse si el modelo mantiene su desempeño bajo el
protocolo cronológico. Y un buen valor global de área bajo la curva no
implica que la probabilidad esté calibrada, ni que el modelo se comporte de
igual manera en las etapas tempranas y tardías de la partida.

\razon{Este apartado existe para impedir que los resultados se empleen fuera
del ámbito en el cual fueron validados. Su omisión trasladaría al lector la
responsabilidad de inferir los límites.}

\subsection{Riesgos de calidad no aislables}
\label{subsec:calidadnoaislable}

La auditoría identificó 73 episodios con desenlace anómalo, cinco episodios
declarados en un manifiesto pero ausentes de los ficheros, y episodios de
duración extrema ---hasta 1\,035 turnos--- cuya causa no puede determinarse
sin inspección individual. Estos casos se documentan y se reservan para
análisis de sensibilidad, pero no puede descartarse la presencia de
comportamientos degenerados de los agentes que no resulten identificables por
completo.

Xenopoulos \textit{et al.} reconocen una limitación análoga al señalar que no
les resulta posible identificar a los jugadores que incurren en trampa
\cite{xenopoulos_winprob_2022}. Bia\l ecki \textit{et al.} declaran que
deliberadamente no filtraron repeticiones terminadas de forma prematura, con
el fin de preservar las distribuciones originales
\cite{bialecki_sc2egset_2023}.

\razon{Declarar las anomalías que se sabe que existen pero no se pueden
aislar por completo es preferible a presentar el corpus como libre de
contaminación. Los precedentes citados establecen que esta declaración es
práctica aceptada en el dominio.}

\subsection{Mitigaciones previstas antes del entrenamiento}
\label{subsec:mitigacionetica}

\begin{itemize}
  \item Toda cifra reportada lleva su alcance: los diez días auditados, la
        franja de puntuación y las versiones del motor presentes.
  \item La discriminación y la calibración se reportan por separado, y los
        resultados bajo los protocolos temporal y de reserva de composiciones
        se reportan aparte del protocolo primario en lugar de promediarse con
        él.
  \item Los identificadores de agente y equipo quedan excluidos del vector de
        características, y no se publican resultados por competidor.
  \item Los intervalos de confianza se calculan sobre episodios, de modo que
        la incertidumbre no se subestime contando filas correlacionadas como
        evidencia independiente.
  \item El resumen, las conclusiones y cualquier síntesis pública declaran la
        población sobre la que se ajustó el modelo, de modo que la limitación
        viaje con el resultado en lugar de quedar confinada a esta sección.
\end{itemize}

\razon{Las cinco mitigaciones son verificables en el propio artículo. Una
mitigación cuya ejecución no puede comprobarse no constituye una mitigación.}

\subsection{Aspectos diferidos}
\label{subsec:diferido}

La ética del despliegue, los escenarios de uso indebido y los riesgos
operativos de un agente que consuma este estimador ---en particular el efecto
de un evaluador sobreconfiado sobre las decisiones del agente--- se abordan en
la etapa final, cuando exista un agente.